% Part: computability
% Chapter: computability-theory
% Section: k-1

\documentclass[../../../include/open-logic-section]{subfiles}

\begin{document}

\olfileid[tr]{cmp}{thy}{k1}
\olsection{Bir İndirgenebilirlik Örneği}

\olref[ppr]{prop:reduce} önermesinin bir uygulamasını ele alalım.

\begin{prop}
\ollabel{prop:k1}
\[
  K_1=\Setabs{e}{\cfind{e}(0)\fdefined}
\]
olsun. O zaman $K_1$ hesaplanabilir biçimde numaralandırılabilir, fakat
hesaplanabilir değildir.
\end{prop}

\begin{proof}
$K_1=\Setabs{e}{\lexists[s][T(e,0,s)]}$ olduğundan
\olref[eqc]{thm:exists-char} uyarınca $K_1$ hesaplanabilir biçimde
numaralandırılabilir.

$K_1$'in hesaplanabilir olmadığını göstermek için $K_0$'ın ona
indirgenebilir olduğunu gösterelim.

\begin{explain}
Bu biraz güçtür; çünkü $K_1$ aracılığıyla yalnızca belirli bir $0$ girdisinde
başlayan hesaplar hakkında soru sorabiliriz. Bu tür soruları yanıtlayabilen
zeki bir arkadaşınız olduğunu varsayın; böyle arkadaşlara ``kâhin'' denir.
Birisi gelip $\tuple{e,x}$ çiftinin $K_0$ içinde olup olmadığını, yani
$e$ makinesinin $x$ girdisinde durup durmadığını sorsun. Yapabileceğiniz şey,
\emph{herhangi bir} girdiyi yok sayıp onun yerine $e$'yi $x$ girdisinde
çalıştıran başka bir $e_x$ makinesi kurmaktır. O zaman $e$ makinesinin $x$
girdisinde durup durmadığı sorusu, $e_x$ makinesinin $0$ girdisinde ya da
başka herhangi bir girdide durup durmadığı sorusuna eşdeğerdir. Arkadaşınıza
bu yeni $e_x$ makinesinin $0$ girdisinde durup durmadığını sorarsınız;
değiştirilmiş soruya verdiği yanıt ilk soruyu yanıtlar. Bu, $K_0$'ın $K_1$'e
istenen indirgemesini sağlar.
\end{explain}

Evrensel kısmi hesaplanabilir fonksiyonu kullanarak $f$'yi
\[
  f(x,y,z)\simeq\cfind{x}(y)
\]
ile tanımlanan üç yerli fonksiyon yapalım. $f$'nin üçüncü girdisini bütünüyle
yok saydığına dikkat edin. $f=\cfind{e}[3]$ olacak bir $e$ indisi seçelim:
\[
  \cfind{e}[3](x,y,z)\simeq\cfind{x}(y).
\]
$s$-$m$-$n$ teoremi uyarınca her $z$ için
\begin{align*}
  \cfind{s(e,x,y)}(z) & \simeq\cfind{e}[3](x,y,z)\\
  & \simeq\cfind{x}(y)
\end{align*}
olacak bir $s(e,x,y)$ fonksiyonu vardır.

\begin{explain}
Yukarıdaki biçimsel olmayan savın terimleriyle $s(e,x,y)$, herhangi bir
$z$ girdisini yok sayıp $\cfind{x}(y)$ değerini hesaplayan makinenin indisidir.
\end{explain}

Özellikle
\[
  \cfind{s(e,x,y)}(0)\fdefined\quad\text{ancak ve ancak}\quad
  \cfind{x}(y)\fdefined
\]
olur. Başka bir deyişle $\tuple{x,y}\in K_0$ ancak ve ancak
$s(e,x,y)\in K_1$ olur. Dolayısıyla
\[
  g(w)=s(e,(w)_0,(w)_1)
\]
ile tanımlanan $g$, $K_0$'ın $K_1$'e bir indirgemesidir.
\end{proof}

\end{document}

